\documentclass[11pt]{article}

% --- Configuración básica ---
\usepackage[spanish,es-nodecimaldot]{babel}
\usepackage[utf8]{inputenc} % si usas LuaLaTeX/XeLaTeX, puedes omitir inputenc
\usepackage[T1]{fontenc}
\usepackage{geometry}
\geometry{margin=2.5cm}
\usepackage{hyperref}
\hypersetup{colorlinks=true, linkcolor=black, urlcolor=blue}
\usepackage{enumitem}
\usepackage{verbatim}
\setlist{itemsep=2pt, topsep=4pt}

\title{Guía breve para usar \texttt{FdC\_notebooks} en JupyterHub (MACTI, UNAM)}
\author{}
\date{}

\begin{document}
\maketitle

Esta guía asume que trabajas en JupyterHub de MACTI, que abre un contenedor Linux ya configurado con los paquetes necesarios para estos notebooks (no necesitas instalar nada adicional).

\begin{enumerate}
  \item \textbf{Inicia sesión en JupyterHub (MACTI)} y arranca tu servidor. Verás el explorador de archivos de tu contenedor.

  \item \textbf{Abre una Terminal} desde JupyterHub:
  \begin{quote}
  Menú \texttt{New} \textrightarrow{} \texttt{Terminal}.
  \end{quote}

  \item \textbf{Clona el repositorio y entra a la carpeta}:
\begin{verbatim}
git clone https://github.com/JLGarciaFranco/FdC_notebooks/tree/main
cd FdC_notebooks
\end{verbatim}


  \item \textbf{Abre los notebooks} desde el panel de JupyterHub: navega a \texttt{FdC\_notebooks/} y haz clic en el \texttt{.ipynb} que quieras (por ejemplo, \texttt{Insolacion.ipynb}, \texttt{Bolton\_approx.ipynb}, \texttt{CAPE.ipynb}).

  \item \textbf{Ejecuta las celdas} normalmente:
  \begin{itemize}
    \item \texttt{Shift+Enter} para ejecutar una celda y avanzar.
    \item Menú \texttt{Run} \textrightarrow{} \texttt{Run All Cells} para correr todo el notebook.
  \end{itemize}

  \item \textbf{Guarda tu trabajo} con el ícono \texttt{Save} o \texttt{Cmd/Ctrl+S}. (Los paquetes requeridos ya vienen instalados en el contenedor de MACTI; no necesitas pasos extra.)

  \item \textbf{Actualizar los notebooks cuando el repositorio cambie} (\emph{pull}):
  \begin{enumerate}[label*=\alph*)]
    \item Cierra los notebooks abiertos y abre una Terminal dentro de \texttt{FdC\_notebooks/}.
    \item Si \emph{no} has modificado archivos del repositorio:
\begin{verbatim}
git pull origin main
\end{verbatim}
    \item Si \emph{sí} hiciste cambios en archivos del repo, pueden surgir conflictos. Para evitarlos:
      \begin{itemize}
        \item Evita editar los notebooks “originales”. Trabaja en copias (ver paso 8).
        \item Si ya editaste un archivo original y aparece un conflicto, respalda tu versión (por ejemplo, renómbralo con \texttt{\_miversion}) y luego ejecuta:
\begin{verbatim}
git restore --staged .
git checkout -- .
git pull origin main
\end{verbatim}
        \item Alternativamente, mueve tus cambios a tu carpeta de trabajo (paso 8) y luego haz \texttt{git pull}.
      \end{itemize}
    \item Tras el \texttt{pull}, recarga el navegador y abre de nuevo los notebooks.
  \end{enumerate}

  \item \textbf{Crea tus propios notebooks de trabajo} para no tener conflictos al actualizar:
  \begin{enumerate}[label*=\alph*)]
    \item Crea una carpeta, por ejemplo \texttt{mi\_trabajo/}, dentro de \texttt{FdC\_notebooks/}.
    \item En el panel de JupyterHub: \texttt{New} \textrightarrow{} \texttt{Notebook} (\texttt{Python 3}).
    \item Trabaja y guarda tus cambios sólo dentro de \texttt{mi\_trabajo/}. Así podrás hacer \texttt{git pull} sin sobrescribir tu trabajo.
  \end{enumerate}
\end{enumerate}

\paragraph{Notas.}
(1) Si \texttt{git pull} indica que no encuentra la rama \texttt{main}, verifica el nombre de la rama por defecto del repo (a veces es \texttt{master}). (2) Si prefieres traer un notebook específico sin mezclar todo, puedes descargarlo directamente desde GitHub y colocarlo en tu carpeta \texttt{mi\_trabajo/}.

\end{document}
